\section{Solutions for Chapter~\ref{ch:branching}}

\paragraph{Exercise~\ref{ex:c-branch-1}: the $x_4$ top.}
\phantomsection
\label{sol:ex:c-branch-1}
At $m=6$, the lattice--Potts contribution is $36/156=3/13$.
The paired parafermion charge is $k=-5\cdot6=-4$ modulo thirteen,
and its energy is
$4\cdot5/15-16/13=4/3-16/13=(52-48)/39=4/39$. The sum is
$9/39+4/39=1/3$. The opposite charge supplies the conjugate
minimizer. Since these are distinct lattice charge sectors and each
has one top state here, their contributions add to a two-dimensional
top. They are not two copies produced by left/right bulk pairing.



\paragraph{Exercise~\ref{ex:c-branch-2}: checking the top certificate.}
\phantomsection
\label{sol:ex:c-branch-2}
The coefficient $4680(17/10)$ equals $7956$. For the requested
three variables the left-side coefficients are
\[
 x_0:7956,\qquad
 x_5:14(7956)-148560=-37176,\qquad
 y_4:14(7956)-81480=29904.
\]
The right side gives respectively $7956$,
$-8460-28716=-37176$, and $28716+1188=29904$.
In particular the negative $x_5$ coefficient is absorbed by
the actual representation constraints
$y_2-x_5\geq0$ and $y_4-x_5-x_6\geq0$.

Ignoring those constraints and taking a weighted average of
individual row ratios gives only the bound by their maximum.
The $x_5$ row alone has ratio $619/273>2$. Such a bound could
not exclude an additional integer-graded module beginning at
weight two. The current character inequalities force the high
Cartan charges of that row to come with lower-charge companions;
their combined contribution yields the stronger bound $17/10$.
The individual ratio is bookkeeping, not a separate ambient
energy eigenvalue.



\paragraph{Exercise~\ref{ex:c-branch-3}: multiplicities are squared.}
\phantomsection
\label{sol:ex:c-branch-3}
Schur's lemma gives $\Hom(X_i,X_j)=0$ for $i\ne j$ and
$\End(X_i)=\C$. Hence the endomorphism algebra is
$\bigoplus_iM_{m_i}(\C)$ and has dimension $\sum_i m_i^2$.
The branching $1+X+Y+Y^*$ therefore gives endomorphism dimension
$1+1+1+1=4$ if the four simples differ, and dimension
$1+1+2^2=6$ when it becomes $1+X+2Y$.
The boundary-rank formula established in the chapter equates these
dimensions with the two possible boundary ranks. This latter
identification is an extension-theory result, not part of Schur's
lemma for an arbitrary semisimple object.



\paragraph{Exercise~\ref{ex:c-branch-4}: the complex fusion algebra.}
\phantomsection
\label{sol:ex:c-branch-4}
In a one-dimensional representation, write the images as scalars
$a,r$. The quadratic relation rules out $r=0$, since that would
give $0=1$. The cross relation now gives $ra=a^2r$, hence
$a=a^2$; together with $a^3=1$, this forces $a=1$.
The remaining quadratic is $r^2-3mr-1=0$, whose roots are
$d_m$ and $-d_m^{-1}$. For the two-dimensional representation take
\[
 \alpha=\begin{pmatrix}\zeta_3&0\\0&\zeta_3^2\end{pmatrix},
 \qquad\rho=\begin{pmatrix}0&1\\1&0\end{pmatrix}.
\]
The cross relation is immediate, and
$(1+\alpha+\alpha^2)\rho=0$, so the quadratic relation reduces to
$\rho^2=1$. Also $\alpha^3=1$ entry by entry. The two distinct
eigenlines of $\alpha$ are exchanged by $\rho$, so neither is a
common invariant line; the representation is irreducible.
These inequivalent irreducibles have squared dimensions
$1+1+4=6$, exhausting the six-dimensional semisimple complex fusion
algebra.



\paragraph{Exercise~\ref{ex:c-branch-5}: the arithmetic exclusion.}
\phantomsection
\label{sol:ex:c-branch-5}
Expanding the difference gives
\begin{align*}
 36m^2(9m^2+4)-13(4m^2+2)^2
 &=116m^4-64m^2-52\\
 &=4(m^2-1)(29m^2+13).
\end{align*}
For $m>1$ it is strictly positive. Both sides of the preceding
root-sum inequality are positive, so squaring would give the
opposite weak inequality
$13(4m^2+2)^2\geq36m^2(9m^2+4)$; this is the contradiction outside
the exceptional squareclass.

For $m=2$, $\sqrt{9m^2+4}=2\sqrt{10}$, and the second moment is
\[
 \frac{-4+(2\varepsilon-2)2\sqrt{10}}2
 =\begin{cases}
 -2,&\varepsilon=1,\\
 -2-4\sqrt{10},&\varepsilon=-1.
 \end{cases}
\]
The first sum is $-2-2\sqrt{10}$, expressed with $4m^2+2=18$
roots. The conductor-forty bound $a+8b$ is exactly
$2+8(2)=18$, so its equality condition requires the squared-root
sum to be $a=2$. The second-moment formulas do lie in
$\Q(\sqrt{10})$, as required for that equality condition.
Neither permitted indicator sign gives the required value.
Thus the exceptional case does not survive.



\paragraph{Exercise~\ref{ex:c-branch-6}: center global dimension.}
\phantomsection
\label{sol:ex:c-branch-6}
The sum is
\[
 1+d^4+4(1+d^2)^2+6(3d)^2.
\]
Squaring $d^2-1=3d$ gives $d^4-2d^2+1=9d^2$, hence
$d^4-11d^2+1=0$. Subtracting
the displayed sum from $9(1+d^2)^2$ gives
$4(d^4-11d^2+1)=0$. Thus the center global dimension is the square
of the boundary global dimension $3(1+d^2)$. Its total quantum
dimension is the positive square root, namely $3(1+d^2)$. Its
rank is twelve; these are three different invariants.


